\documentclass[12pt]{article}

\usepackage[a4paper,margin=1in]{geometry}
\usepackage{amsmath,amssymb}
\usepackage{newtxtext,newtxmath} % Times-style font
\usepackage{fancyhdr}
\usepackage{listings}
\usepackage{graphicx}
\usepackage{float}
\usepackage{booktabs}
\usepackage{xcolor}

\setlength{\parindent}{0pt}
\setlength{\parskip}{5pt}

% Keep entire document at 12 pt
\AtBeginDocument{
    \fontsize{12pt}{14pt}\selectfont
}

% ------------------------------------------------
% Header
% ------------------------------------------------

\pagestyle{fancy}
\fancyhf{}

\fancyhead[L]{Name: \underline{Shaik Jnaved}}
\fancyhead[R]{Enrollment No.: \underline{P26DS001}}

\renewcommand{\headrulewidth}{0.4pt}
\setlength{\headheight}{15pt}

% ------------------------------------------------
% Code formatting
% ------------------------------------------------

\lstset{
    language=C,
    basicstyle=\ttfamily\fontsize{9pt}{11pt}\selectfont,
    breaklines=true,
    frame=single,
    columns=fullflexible,
    keepspaces=true,
    showstringspaces=false,
    keywordstyle=\textbf,
    tabsize=4,
    numbers=left,
    numberstyle=\tiny\color{gray},
    numbersep=8pt
}

\begin{document}

% ==============================================================================
% TITLE
% ==============================================================================

\begin{center}
\textbf{Computer Science and Engineering Department, SVNIT, Surat}\\
\textbf{M.Tech. I -- Semester I}\\
\textbf{High Performance Computing}\\
\textbf{Lab Assignment 3}
\end{center}

\vspace{0.4cm}

\textbf{Problem Statement}

\textbf{Title: Parallel 3D Linear Regression using OpenMP and Sequential Performance Comparison}

Implement a parallel version of multiple linear regression using OpenMP on a 3-dimensional dataspace ($X_1, X_2 \in \mathbb{R}$ features, and $Y \in \mathbb{R}$ label). Apply data-parallel gradient descent to optimize parameters $w_1, w_2, b$ minimizing the Mean Squared Error (MSE) loss function over $N = 2,000,000$ synthetic datapoints and $100$ epochs. Evaluate and compare execution runtime, speedup, and parallel efficiency across varying OpenMP thread configurations ($1, 2, 4, 8$ threads) against the baseline sequential algorithm.

\vspace{0.3cm}

\textbf{Source Code:}

\begin{lstlisting}
/*
 * High Performance Computing - Lab Assignment 3
 * Question 1: Parallel Linear Regression using OpenMP and Sequential Comparison
 * 
 * Problem Statement:
 * Implement a parallel version of linear regression using OpenMP.
 * Consider X values are two-dimensional and Y label is one-dimensional,
 * i.e., overall dataspace is 3-dimensional (X1, X2, Y).
 * Apply parallelization on Gradient Descent computation and compare its
 * running time performance using OpenMP with its sequential counterpart.
 * Use a significantly large number of datapoints for performance comparison.
 * 
 * Compilation & Execution:
 *   gcc -O3 -fopenmp lab_3/q1.c -o lab_3/q1.exe -lm
 *   .\lab_3\q1.exe
 */

#include <stdio.h>
#include <stdlib.h>
#include <math.h>
#include <time.h>
#include <omp.h>

// ==========================================
// CONFIGURATION PARAMETERS
// ==========================================
#define N 2000000              // Number of data points (2 Million)
#define NUM_EPOCHS 100         // Number of gradient descent iterations
#define LEARNING_RATE 0.1f     // Learning rate (alpha / eta)
#define NUM_THREADS 8          // Number of OpenMP threads for parallel execution

// Ground truth parameters for synthetic dataset: Y = W1_TRUE * X1 + W2_TRUE * X2 + B_TRUE + noise
#define W1_TRUE 3.50f
#define W2_TRUE 1.50f
#define B_TRUE  2.00f

// ==========================================
// DATA GENERATION & LOSS FUNCTION
// ==========================================
void generate_dataset(float *x1, float *x2, float *y, int n) {
    srand(42); // Fixed seed for reproducible benchmarks
    for (int i = 0; i < n; i++) {
        // Normalized feature coordinates in [0, 1]
        x1[i] = (float)rand() / (float)RAND_MAX;
        x2[i] = (float)rand() / (float)RAND_MAX;

        // Small noise in [-0.05, 0.05]
        float noise = (((float)rand() / (float)RAND_MAX) - 0.5f) * 0.1f;
        y[i] = (W1_TRUE * x1[i]) + (W2_TRUE * x2[i]) + B_TRUE + noise;
    }
}

// Compute Mean Squared Error (MSE) loss: J(w1, w2, b) = 1/(2N) * sum((y_pred - y_actual)^2)
float compute_mse_loss(const float *x1, const float *x2, const float *y, int n, float w1, float w2, float b) {
    double total_error = 0.0;
    #pragma omp parallel for reduction(+:total_error)
    for (int i = 0; i < n; i++) {
        float y_pred = (w1 * x1[i]) + (w2 * x2[i]) + b;
        float err = y_pred - y[i];
        total_error += (double)(err * err);
    }
    return (float)(total_error / (2.0 * (double)n));
}

// ==========================================
// 1. SEQUENTIAL GRADIENT DESCENT
// ==========================================
double linear_regression_sequential(const float *x1, const float *x2, const float *y, int n, int epochs, float alpha, float *final_w1, float *final_w2, float *final_b) {
    float w1 = 0.0f;
    float w2 = 0.0f;
    float b  = 0.0f;

    double start_time = omp_get_wtime();

    for (int epoch = 0; epoch < epochs; epoch++) {
        float dw1 = 0.0f;
        float dw2 = 0.0f;
        float db  = 0.0f;

        // Sequential accumulation of partial gradients over N datapoints
        for (int i = 0; i < n; i++) {
            float y_pred = (w1 * x1[i]) + (w2 * x2[i]) + b;
            float err = y_pred - y[i];

            dw1 += err * x1[i];
            dw2 += err * x2[i];
            db  += err;
        }

        // Parameter update step
        w1 -= alpha * (dw1 / (float)n);
        w2 -= alpha * (dw2 / (float)n);
        b  -= alpha * (db  / (float)n);
    }

    double end_time = omp_get_wtime();

    *final_w1 = w1;
    *final_w2 = w2;
    *final_b  = b;

    return (end_time - start_time);
}

// ==========================================
// 2. OPENMP PARALLEL GRADIENT DESCENT
// ==========================================
double linear_regression_openmp(const float *x1, const float *x2, const float *y, int n, int epochs, float alpha, int num_threads, float *final_w1, float *final_w2, float *final_b) {
    float w1 = 0.0f;
    float w2 = 0.0f;
    float b  = 0.0f;

    double start_time = omp_get_wtime();

    for (int epoch = 0; epoch < epochs; epoch++) {
        float dw1 = 0.0f;
        float dw2 = 0.0f;
        float db  = 0.0f;

        // Parallel data decomposition with static schedule and reduction on gradient accumulators
        #pragma omp parallel for reduction(+:dw1, dw2, db) num_threads(num_threads) schedule(static)
        for (int i = 0; i < n; i++) {
            float y_pred = (w1 * x1[i]) + (w2 * x2[i]) + b;
            float err = y_pred - y[i];

            dw1 += err * x1[i];
            dw2 += err * x2[i];
            db  += err;
        }

        // Parameter update step
        w1 -= alpha * (dw1 / (float)n);
        w2 -= alpha * (dw2 / (float)n);
        b  -= alpha * (db  / (float)n);
    }

    double end_time = omp_get_wtime();

    *final_w1 = w1;
    *final_w2 = w2;
    *final_b  = b;

    return (end_time - start_time);
}

// ==========================================
// MAIN FUNCTION
// ==========================================
int main() {
    printf("=========================================================================\n");
    printf("     High Performance Computing - Linear Regression Optimization        \n");
    printf("=========================================================================\n");
    printf("Configuration:\n");
    printf("  Dataspace Dimension : 3D (2D Features: X1, X2 | 1D Target: Y)\n");
    printf("  Number of Points (N): %d\n", N);
    printf("  Number of Epochs    : %d\n", NUM_EPOCHS);
    printf("  Learning Rate (eta) : %.4f\n", LEARNING_RATE);
    printf("  OpenMP Threads      : %d\n", NUM_THREADS);
    printf("  True Model Equation : Y = %.2f * X1 + %.2f * X2 + %.2f\n", W1_TRUE, W2_TRUE, B_TRUE);
    printf("-------------------------------------------------------------------------\n");

    // Dynamic memory allocation on host
    float *h_x1 = (float *)malloc(N * sizeof(float));
    float *h_x2 = (float *)malloc(N * sizeof(float));
    float *h_y  = (float *)malloc(N * sizeof(float));

    if (!h_x1 || !h_x2 || !h_y) {
        printf("Error: Failed to allocate memory for dataset.\n");
        return -1;
    }

    printf("Generating synthetic 3D dataset...\n");
    generate_dataset(h_x1, h_x2, h_y, N);
    printf("Dataset generated successfully.\n\n");

    // ---------------------------------------------------------
    // 1. Sequential Execution
    // ---------------------------------------------------------
    printf("[1] Running Sequential Linear Regression...\n");
    float seq_w1, seq_w2, seq_b;
    double seq_time = linear_regression_sequential(h_x1, h_x2, h_y, N, NUM_EPOCHS, LEARNING_RATE, &seq_w1, &seq_w2, &seq_b);
    float seq_loss = compute_mse_loss(h_x1, h_x2, h_y, N, seq_w1, seq_w2, seq_b);
    printf("    Learned Model : Y = %.4f * X1 + %.4f * X2 + %.4f\n", seq_w1, seq_w2, seq_b);
    printf("    Final MSE Loss: %.6f\n", seq_loss);
    printf("    Execution Time: %.6f seconds\n\n", seq_time);

    // ---------------------------------------------------------
    // 2. OpenMP Parallel Execution
    // ---------------------------------------------------------
    printf("[2] Running OpenMP Parallel Linear Regression (%d Threads)...\n", NUM_THREADS);
    float omp_w1, omp_w2, omp_b;
    double omp_time = linear_regression_openmp(h_x1, h_x2, h_y, N, NUM_EPOCHS, LEARNING_RATE, NUM_THREADS, &omp_w1, &omp_w2, &omp_b);
    float omp_loss = compute_mse_loss(h_x1, h_x2, h_y, N, omp_w1, omp_w2, omp_b);
    printf("    Learned Model : Y = %.4f * X1 + %.4f * X2 + %.4f\n", omp_w1, omp_w2, omp_b);
    printf("    Final MSE Loss: %.6f\n", omp_loss);
    printf("    Execution Time: %.6f seconds\n", omp_time);
    printf("    Speedup vs Seq: %.2fx\n\n", seq_time / omp_time);

    // ---------------------------------------------------------
    // 3. Granularity Analysis (Thread Scaling)
    // ---------------------------------------------------------
    printf("-------------------------------------------------------------------------\n");
    printf("                    THREAD SCALABILITY COMPARISON                        \n");
    printf("-------------------------------------------------------------------------\n");
    printf("%-15s | %-12s | %-12s | %-12s\n", "Thread Count", "Time (s)", "Speedup", "Efficiency");
    printf("-------------------------------------------------------------------------\n");
    printf("%-15s | %10.6f s | %10.2fx | %10.1f%%\n", "1 (Sequential)", seq_time, 1.0, 100.0);

    int thread_counts[] = {2, 4, 8};
    int num_configs = sizeof(thread_counts) / sizeof(thread_counts[0]);

    for (int k = 0; k < num_configs; k++) {
        int threads = thread_counts[k];
        float t_w1, t_w2, t_b;
        double t_time = linear_regression_openmp(h_x1, h_x2, h_y, N, NUM_EPOCHS, LEARNING_RATE, threads, &t_w1, &t_w2, &t_b);
        double speedup = seq_time / t_time;
        double efficiency = (speedup / (double)threads) * 100.0;
        
        char label[20];
        snprintf(label, sizeof(label), "%d Threads", threads);
        printf("%-15s | %10.6f s | %10.2fx | %10.1f%%\n", label, t_time, speedup, efficiency);
    }
    printf("=========================================================================\n");

    // Clean up allocated memory
    free(h_x1);
    free(h_x2);
    free(h_y);

    return 0;
}
\end{lstlisting}

\newpage

\textbf{Output:}

\begin{figure}[H]
\centering
\includegraphics[width=0.95\textwidth]{images/q1.png}
\caption{Execution and Scalability Benchmark for Parallel Linear Regression ($N = 2,000,000$, Epochs = $100$)}
\end{figure}

\vspace{0.3cm}

\textbf{Analysis:}

\textbf{1. OpenMP Parallel Data Decomposition \& Reduction:}
\begin{itemize}
    \item \textbf{Work-Sharing Directive:} \texttt{\#pragma omp parallel for reduction(+:dw1, dw2, db) num\_threads(T) schedule(static)} parallelizes the inner loop computing error contributions across $N = 2,000,000$ points.
    \item \textbf{Static Scheduling \& Cache Locality:} With \texttt{schedule(static)}, iterations are partitioned into contiguous blocks of size $\lfloor N / T \rfloor$. Because arrays \texttt{x1}, \texttt{x2}, and \texttt{y} are stored in contiguous 1D memory buffers, sequential thread accesses maximize L1/L2 hardware prefetching and eliminate runtime task scheduling overhead.
    \item \textbf{Multi-Variable Reduction:} The \texttt{reduction(+:dw1, dw2, db)} clause creates private copies of \texttt{dw1}, \texttt{dw2}, and \texttt{db} for each thread, initialized to \texttt{0.0f}. Each thread accumulates partial sums locally in CPU registers/private cache lines. At loop termination, OpenMP combines all private copies into the master variables using a lock-free tree reduction, avoiding race conditions and eliminating cache-line false sharing.
\end{itemize}

\textbf{2. Scalability, Speedup, and Efficiency Analysis:}
Parallel performance metrics are evaluated using:
\[
\text{Speedup } S_p = \frac{T_1}{T_p}, \qquad \text{Parallel Efficiency } E_p = \frac{S_p}{p} \times 100\%
\]
\begin{itemize}
    \item \textbf{2 Threads:} Execution time decreases from $0.161999\text{ s}$ to $0.084922\text{ s}$, achieving a speedup of $\mathbf{1.91\times}$ and near-ideal parallel efficiency of $\mathbf{95.4\%}$.
    \item \textbf{4 Threads:} Execution time drops to $0.049519\text{ s}$, yielding a speedup of $\mathbf{3.27\times}$ and an efficiency of $\mathbf{81.8\%}$.
    \item \textbf{8 Threads:} Execution time reaches $0.048290\text{ s}$ ($\mathbf{3.35\times}$ speedup, $\mathbf{41.9\%}$ efficiency).
\end{itemize}

\textbf{3. Architectural Bottlenecks and Saturation (Memory Bandwidth \& Amdahl's Law):}
\begin{itemize}
    \item \textbf{Memory-Bound Nature:} Each epoch streams $3 \times 2,000,000 \times 4\text{ bytes} = 24\text{ MB}$ of data across the memory bus. Across $100$ epochs, $2.4\text{ GB}$ is streamed. At 8 threads, the arithmetic intensity is low ($\approx 5$ FLOPs per 12 bytes read = $0.42\text{ FLOP/byte}$), leading to memory bus bandwidth saturation where ALU cores stall waiting for RAM/LLC data prefetching.
    \item \textbf{Inter-Epoch Barrier Synchronization:} At the end of each epoch, all threads synchronize at the reduction barrier while the master thread executes the scalar parameter update. This sequential fraction imposes an asymptotic speedup ceiling as dictated by Amdahl's Law.
\end{itemize}

\begin{table}[H]
\centering
\begin{tabular}{lccccc}
\toprule
\textbf{Configuration} & \textbf{Threads ($p$)} & \textbf{Execution Time ($T_p$)} & \textbf{Speedup ($S_p$)} & \textbf{Efficiency ($E_p$)} & \textbf{Final MSE Loss} \\
\midrule
Sequential Baseline & 1 & $0.161999\text{ s}$ & $1.00\times$ & $100.0\%$ & $0.040441$ \\
OpenMP Parallel     & 2 & $0.084922\text{ s}$ & $1.91\times$ & $95.4\%$  & $0.040423$ \\
OpenMP Parallel     & 4 & $0.049519\text{ s}$ & $3.27\times$ & $81.8\%$  & $0.040423$ \\
OpenMP Parallel     & 8 & $0.048290\text{ s}$ & $3.35\times$ & $41.9\%$  & $0.040423$ \\
\bottomrule
\end{tabular}
\caption{Performance and Thread Scalability Summary on 2,000,000 Datapoints ($100$ Epochs)}
\end{table}

\end{document}
